%=====================================================================
\section{Alignment and Prosody Marking}
%=====================================================================

This stage recovers the speaker's delivery from the audio and writes it back into the Bengali text as inline markers. It first locates each word in time by forced alignment, and it then analyzes the word timings and the raw waveform with fixed rules to mark pauses and to measure loudness and elongation. Only the alignment step uses a model.

\subsection{Forced alignment}
Forced alignment determines where each known word occurs in the audio. Because the aligner (Table~\ref{tab:models}) was trained on romanized text in many languages, it can be used for Bengali without a language-specific aligner.

The audio is mixed to mono and resampled to \textbf{16\,kHz}. The model produces character probabilities for every short frame of audio, and the aligner finds the most likely path of the known characters through those frames. This gives each word a start time and an end time, which the following steps use.

\subsection{Pauses}
The silence before each word is measured and compared with two fixed thresholds:
\begin{itemize}[nosep]
  \item A \textbf{short pause}, marked with a comma (\texttt{,}), is any gap of at least \textbf{0.20\,s} and shorter than 0.50\,s.
  \item A \textbf{long pause}, marked with a bar (\texttt{|}), is any gap of at least \textbf{0.50\,s}.
\end{itemize}
The silence before the first word and the silence after the last word are not marked.

\subsection{Loudness}
The waveform is split into \textbf{25\,ms} frames, and the energy of each frame is measured in decibels. A word's loudness is the average of its loudest frames, so that brief dips do not pull it down. Each word is then compared with the rest of the clip. Words well above the clip's typical loudness are marked \emph{loud}, and words well below it are marked \emph{soft}.

\subsection{Elongated words}
Speakers often elongate a word for effect. Since a longer word naturally takes longer to say, its duration is divided by its romanized length to obtain a time per character. Each word is then compared with the clip's median time per character, and a word that is unusually slow for its length is marked \emph{dragged}.

\subsection{The marked sentence}
The Bengali words are rebuilt into a sentence, with each pause marker placed before the word that follows the pause. \textbf{In the current procedure only the pause markers reach translation and synthesis.} Loudness and elongation are still measured and reviewed.

\begin{example}[: alignment and prosody marking]
The first diagram shows the result of the stage on the time axis of the clip. Each word has its own column, which follows a dotted guide from its aligned position down to its loudness relative to the clip average and its duration relative to the duration expected for its romanized length.

\begin{center}
% Worked example, alignment and prosody marking: one panel on the time axis of the clip
% (1 s = 1.45 cm), with one column per word. Each column holds the word's aligned block (below
% the audio envelope), its name, its loudness bar and its length bar, joined by a faint vertical
% guide. Names are vertical (90 degrees) and small, hanging below their blocks. Row titles are on
% the left and the legend is on the right. Row positions (y, cm):
%   audio envelope 0, word blocks -0.55, names from -0.80 down,
%   loudness baseline -3.00, length baseline -5.70, time axis -5.85.
\begin{tikzpicture}[x=1.45cm,
    emph none/.style={fill=black!25},
    emph loud/.style={fill=red!60},
    emph soft/.style={fill=blue!50},
    drag no/.style={fill=black!25},
    drag yes/.style={fill=orange!75},
    flag none/.style={draw=black!50},
    flag gap/.style={draw=red!75!black, line width=1pt},
    flag short/.style={flag gap},
    label none/.style={},
    label gap/.style={text=red!75!black},
    label short/.style={label gap},
    panel/.style={fill=black!3, draw=black!15, rounded corners=3pt},
    row title/.style={font=\footnotesize\bfseries, align=right, anchor=east},
    legend/.style={font=\scriptsize, anchor=west},
    guide/.style={black!22, densely dotted},
  ]
  \pgfmathtruncatemacro\lastsecond{\clipduration}
  \def\legendx{6.75}   % where the legend starts (s)
  \def\wordy{-0.55}    % centre of the word blocks
  \def\loudy{-3.00}    % loudness baseline (clip average)
  \def\lengthy{-5.70}  % length baseline (zero length)

  % One panel for all rows
  \fill[panel] (-1.55,0.55) rectangle (8.95,-6.45);

  % Row titles
  \node[row title] at (-0.1,-0.55)   {Words and\\pauses};
  \node[row title] at (-0.1,\loudy)  {Loudness};
  \node[row title] at (-0.1,-5.05)   {Length vs.\\expected};

  % Faint row separators
  \draw[black!12] (-1.4,-2.15) -- (8.8,-2.15);
  \draw[black!12] (-1.4,-4.10) -- (8.8,-4.10);

  % Guides: one per word, from its block down through the loudness and length rows
  \foreach \word/\tstart/\tend/\flag/\loudz/\emphasis/\ratio/\dragged in \words
    \draw[guide] ({(\tstart+\tend)/2},{\wordy-0.18}) -- ({(\tstart+\tend)/2},\lengthy);

  % Reference lines
  \draw[black!40] (0,\loudy) -- (\clipduration,\loudy);                       % clip average loudness
  \draw[black!40] (0,\lengthy) -- (\clipduration,\lengthy);                   % zero length
  \draw[black!50, dashed] (0,{\lengthy+0.6}) -- (\clipduration,{\lengthy+0.6}); % expected length

  % Audio envelope of the clip, above the word blocks
  \begin{scope}[yscale=0.3]
    \fill[black!20] plot file {figures/data/original-envelope.dat} -- cycle;
  \end{scope}

  % One column per word; words that follow a silence of 0.5 s or more are outlined in red
  \foreach \word/\tstart/\tend/\flag/\loudz/\emphasis/\ratio/\dragged in \words {
    \draw[flag \flag, fill=black!8] (\tstart,{\wordy-0.18}) rectangle (\tend,{\wordy+0.18});
    % the vertical name hangs below the block and is filled with the panel colour,
    % so that it masks the guide behind it
    \node[font=\tiny\ttfamily, rotate=90, anchor=east, inner sep=0.8pt, fill=black!3, label \flag]
      at ({(\tstart+\tend)/2},{\wordy-0.22}) {\word};
    % loudness: bar height follows the z-score (0.3 cm per standard deviation)
    \fill[emph \emphasis] (\tstart,\loudy) rectangle (\tend,{\loudy + 0.3*\loudz});
    % length: 0.6 cm corresponds to the expected length
    \fill[drag \dragged] (\tstart,\lengthy) rectangle (\tend,{\lengthy + 0.6*\ratio});
  }

  % Pause markers: a red bar with its length for a long pause, an orange comma for a short one
  \def\pauselong#1#2{%
    \draw[red!75!black, line width=1.6pt] (#1,{\wordy-0.32}) -- (#1,{\wordy+0.32});
    \node[font=\scriptsize, text=red!75!black, below] at (#1,{\wordy-0.32}) {#2\,s};}
  \def\pauseshort#1#2{%
    \node[font=\Large\bfseries, text=orange!80!black] at (#1,{\wordy-0.12}) {,};}
  \foreach \at/\type/\gap in \pauses {
    \csname pause\type\endcsname{\at}{\gap}
  }

  % Label the word with the largest length ratio
  \gdef\maxratio{0}
  \foreach \word/\tstart/\tend/\flag/\loudz/\emphasis/\ratio/\dragged in \words {
    \pgfmathparse{\ratio>\maxratio ? 1 : 0}
    \ifnum\pgfmathresult=1
      \xdef\maxratio{\ratio}
      \pgfmathsetmacro\middle{(\tstart+\tend)/2}
      \xdef\maxat{\middle}
    \fi
  }
  \node[font=\scriptsize, text=orange!60!black, above] at (\maxat,{\lengthy + 0.6*\maxratio})
    {\pgfmathprintnumber[fixed, precision=1]{\maxratio}×};

  % Time axis, below the length row
  \draw[black!60] (0,-5.85) -- (\clipduration,-5.85);
  \foreach \t in {0,...,\lastsecond}
    \draw[black!60] (\t,-5.85) -- ++(0,-0.08) node[below, font=\scriptsize] {\t\,s};

  % Legend, grouped beside each row
  \begin{scope}[shift={(\legendx,0)}, x=1cm]
    % words and pauses
    \fill[black!20] (0,0.20) rectangle (0.25,0.40);                 \node[legend] at (0.3,0.30) {audio envelope};
    \draw[black!50, fill=black!8] (0,-0.20) rectangle (0.25,0.00); \node[legend] at (0.3,-0.10) {word};
    \draw[red!75!black, line width=1pt, fill=black!8] (0,-0.70) rectangle (0.25,-0.50);
    \node[legend, align=left] at (0.3,-0.60) {word after a silence\\of 0.5\,s or more};
    \draw[red!75!black, line width=1.6pt] (0.12,-1.30) -- (0.12,-1.00); \node[legend] at (0.3,-1.15) {long pause (length)};
    \node[font=\Large\bfseries, text=orange!80!black] at (0.12,-1.65) {,}; \node[legend] at (0.3,-1.60) {short pause};
    % loudness
    \fill[red!60]   (0,-2.60) rectangle (0.25,-2.40); \node[legend] at (0.3,-2.50) {loud};
    \fill[blue!50]  (0,-2.95) rectangle (0.25,-2.75); \node[legend] at (0.3,-2.85) {soft};
    \fill[black!25] (0,-3.30) rectangle (0.25,-3.10); \node[legend] at (0.3,-3.20) {unmarked};
    \draw[black!40] (0,-3.55) -- (0.25,-3.55);        \node[legend] at (0.3,-3.55) {clip average};
    % length
    \fill[orange!75] (0,-4.50) rectangle (0.25,-4.30); \node[legend] at (0.3,-4.40) {dragged};
    \fill[black!25]  (0,-4.85) rectangle (0.25,-4.65); \node[legend] at (0.3,-4.75) {not dragged};
    \draw[black!50, dashed] (0,-5.10) -- (0.25,-5.10); \node[legend] at (0.3,-5.10) {expected length};
  \end{scope}
\end{tikzpicture}

\end{center}

The second diagram compares the silence before each word with the two pause thresholds.

\begin{center}
% Worked example, Stage 3 (prosody marking): the silence before each word against the two pause thresholds.
% The x axis is the length of a silence in seconds (1 s = 8.4 cm, height 1.4 cm); the legend sits on the right.
% Filled dots are gaps between words; hollow dots are the leading and trailing silence.
\begin{tikzpicture}[x=8.4cm, y=1.4cm,
    legend/.style={font=\scriptsize, anchor=west},
  ]
  \def\axisend{1.2}

  % Regions: no pause, short pause and long pause
  \fill[black!6]   (0,0)               rectangle (\shortthreshold,1);
  \fill[orange!15] (\shortthreshold,0) rectangle (\longthreshold,1);
  \fill[red!12]    (\longthreshold,0)  rectangle (\axisend,1);

  % Silences, stacked in four rows so that equal values do not hide each other
  \foreach [count=\i] \gap in \innergaps
    \fill[black!75] (\gap,{0.15 + mod(\i,4)*0.13}) circle (1.6pt);
  \foreach \gap in \edgegaps
    \draw[black!75, fill=white] (\gap,0.35) circle (1.8pt);

  % The two thresholds
  \draw[orange!80!black, line width=1.2pt] (\shortthreshold,0) -- (\shortthreshold,1);
  \draw[red!75!black,    line width=1.2pt] (\longthreshold,0)  -- (\longthreshold,1);

  % Region labels
  \node[font=\scriptsize, text=black!60, fill=black!6, inner sep=1pt]
    at ({\shortthreshold/2},0.85) {no pause};
  \node[font=\scriptsize, text=orange!70!black, fill=orange!15, inner sep=1pt]
    at ({(\shortthreshold+\longthreshold)/2},0.85) {short pause (\texttt{,})};
  \node[font=\scriptsize, text=red!70!black, fill=red!12, inner sep=1pt]
    at ({(\longthreshold+\axisend)/2},0.85) {long pause (\texttt{|})};

  % Axis
  \draw[black!60] (0,0) -- (\axisend,0);
  \foreach \t in {0,0.2,0.4,0.6,0.8,1.0,1.2}
    \draw[black!60] (\t,0) -- ++(0,-0.08) node[below, font=\scriptsize] {\t\,s};

  % Legend, beside the diagram (positions in cm)
  \begin{scope}[shift={(\axisend,0)}, x=1cm, y=1cm, xshift=0.35cm]
    \fill[black!75] (0.1,1.25) circle (1.6pt);
    \node[legend] at (0.25,1.25) {gap between words};
    \draw[black!75, fill=white] (0.1,0.85) circle (1.8pt);
    \node[legend] at (0.25,0.85) {edge silence (not marked)};
    \draw[orange!80!black, line width=1.2pt] (0,0.45) -- (0.2,0.45);
    \node[legend] at (0.25,0.45) {short-pause threshold (\shortthreshold\,s)};
    \draw[red!75!black, line width=1.2pt] (0,0.05) -- (0.2,0.05);
    \node[legend] at (0.25,0.05) {long-pause threshold (\longthreshold\,s)};
  \end{scope}
\end{tikzpicture}

\end{center}

Only the pause marker is written into the transcribed sentence:

\smallskip
\begin{center}{\bn গ্রেগরির কাছ থেকে এখনও উত্তর পেলাম না {\color{red!75!black}\textbf{|}} সে কি তাহলে ইংল্যাণ্ডে নেই}\end{center}
\end{example}
